% --- 1. INFORMAZIONI GENERALI ---
\section{Informazioni Generali}
\begin{itemize}
	\item \textbf{Luogo/Piattaforma:} Server Discord\textsubscript{G}
	\item \textbf{Data:} 2026-05-19
	\item \textbf{Orario di inizio:} 14:30
	\item \textbf{Orario di fine:} 16:40
	\item \textbf{Responsabile\textsubscript{G}} Anton Ricardo Rupi
	\item \textbf{Scriba:} Sofia De Blasi
\end{itemize}

\vspace{0.5cm}
\textbf{Partecipanti presenti:}
\begin{table}[ht]
	\centering
	\renewcommand{\arraystretch}{1.2}
	\begin{tabularx}{\textwidth}{X l}
		\toprule
		\textbf{Nome e Cognome} & \textbf{Matricola} \\
		\midrule
		Sofia De Blasi & 2111014 \\
		Roberto Mariano Doroftei & 2111031 \\
		Filippo Idiometri & 2035617 \\
		Francesco Marcon & 2110990 \\
		Armando Moda Scarati & 2082864 \\
		Anton Ricardo Rupi & 2054304 \\
		\bottomrule
	\end{tabularx}
\end{table}

\vspace{0.5cm}
\textbf{Partecipanti assenti:}
\begin{table}[ht]
	\centering
	\renewcommand{\arraystretch}{1.2}
	\begin{tabularx}{\textwidth}{X l}
		\toprule
		\textbf{Nome e Cognome} & \textbf{Matricola} \\
		\midrule
		Nessuno & - \\
		\bottomrule
	\end{tabularx}
\end{table}

% --- 2. ORDINE DEL GIORNO ---
\section{Ordine del Giorno}
Gli argomenti previsti per la riunione odierna sono:
\begin{enumerate}
	\item Sprint\textsubscript{G} review e retrospective.
	\item Pianificazione dello sprint successivo e assegnazione dei task\textsubscript{G} e ruoli.
\end{enumerate}

% --- 3. RESOCONTO E DISCUSSIONE ---
\section{Resoconto della Riunione}

\subsection{Sprint review\textsubscript{G} e retrospective}
Il gruppo si è confrontato sull’andamento dello
sprint appena concluso. Il lavoro individuale è proceduto linearmente senza riscontrare problemi, 
ma è necessario sviluppare una visione d'insieme del progetto, 
garantendo il pieno allineamento di tutto il team su ogni attività. 
Per quanto riguarda l'ADR, il documento è ufficialmente terminato 
e tutti i membri del team dovranno prenderne visione completa; contestualmente, 
si provvederà a contattare l’azienda per un confronto diretto sulla completezza dei requisiti 
e sulle scelte tecnologiche adottate. Per quanto concerne il Piano di Qualifica 
si avvierà uno studio per capire quali metriche siano richieste, 
distinguendo quelle già in possesso del team per le quali basterà modellare il grafico da quelle mancanti, 
individuando le modalità per ottenere i necessari dati di qualifica e tecnologici. 
Infine lo Stato dell'Arte è stata verificato con successo e approvato.

\subsection{Pianificazione dello sprint successivo e assegnazione dei task e ruoli}
Lo sprint è stato definito con durata dal 19 maggio al 26 maggio. I ruoli assegnati per questo sprint sono: 
Responsabile (Anton Ricardo Rupi), Amministratore (Sofia De Blasi, Filippo Idiometri, Francesco Marcon), Verificatore\textsubscript{G} (Roberto Mariano Doroftei), 
Analisti (Armando Moda Scarati).
Gli amministratori si occuperanno di avviare le comunicazioni con l'azienda proponente per un confronto mirato sull'ADR, sulle scelte tecnologiche e sul perimetro del PoC; 
avranno inoltre il compito di portare avanti la redazione del Piano di Progetto e del Piano di Qualifica. Per quest'ultimo documento, l'attività si concentrerà sull'inserimento 
dei commenti nelle sezioni pertinenti, sulla verifica della completezza delle azioni correttive e sull'analisi delle metriche necessarie per il cruscotto dell'RTB, distinguendo 
quelle già disponibili da quelle mancanti. Infine, gli amministratori provvederanno all'aggiornamento delle Norme di Progetto per quanto concerne l'associazione tra ruoli e documenti. 
L'analista sarà incaricato di ridefinire i Casi d'Uso, spostando UC1 e UC2 (relativi alla creazione e all'eliminazione della nota) all'interno della macroarea 3 come requisiti opzionali; 
tale modifica comporterà la conseguente rinumerazione di tutti gli UC e l'aggiornamento dei relativi riferimenti in ogni sezione del documento, inclusi i target, le tabelle dei requisiti, le fonti e i riepiloghi. 
Parallelamente, tutti i membri del team dedicheranno del tempo allo studio approfondito delle tecnologie selezionate per il PoC. Per concludere, il responsabile si occuperà della stesura del diario di bordo settimanale e della redazione del presente verbale interno.

% --- 4. DECISIONI PRESE ---
\section{Decisioni Prese}

\renewcommand{\arraystretch}{1.3}
\vspace{-0.3cm}
\noindent
\begin{longtable}{c p{12cm}}
	\toprule
	\textbf{Codice} & \textbf{Decisione} \\
	\midrule
	\endfirsthead

	\toprule
	\textbf{Codice} & \textbf{Decisione} \\
	\midrule
	\endhead

	\midrule
	\multicolumn{2}{r}{\textit{Continua nella pagina successiva}} \\
	\midrule
	\endfoot

	\bottomrule
	\endlastfoot

	
    \textbf{VI-14.1} & Scrivere una mail a Zucchetti per parlare di AdR, tecnologie e PoC \\[4pt]

    \textbf{VI-14.2} & Studiare le tecnologie scelte per il PoC \\[4pt]

    \textbf{VI-14.3} & Proseguire il PdQ \\[4pt]

    \textbf{VI-14.4} & Modicare l'AdR \\[4pt]

    \textbf{VI-14.5} & Aggiornare le norme di progetto \\[4pt]

\end{longtable}

% --- 5. AZIONI (TODO) ---
\section{Azioni da Svolgere (To-Do)}

\renewcommand{\arraystretch}{1.3}
\begin{longtable}{c c p{4.5cm} p{3cm} l}
	\toprule
	\textbf{Codice} & \textbf{Rif. Decisione} & \textbf{Descrizione Task} & \textbf{Assegnatario} & \textbf{Scadenza} \\
	\midrule
	\endfirsthead

	\toprule
	\textbf{Codice} & \textbf{Rif. Decisione} & \textbf{Descrizione Task} & \textbf{Assegnatario} & \textbf{Scadenza} \\
	\midrule
	\endhead

	\midrule
	\multicolumn{5}{r}{\textit{Continua nella pagina successiva}} \\
	\midrule
	\endfoot

	\bottomrule
	\endlastfoot

	\ghissue{170} & VI-14.1 & Scrivere email a zucchetti con richiesta riunione in cui parlare di adr, confronto su tecnologie, sul perimetro del poc & Sofia De Blasi & 2026-05-26 \\[4pt]

	\ghissue{171} & - & Verbale interno della riunione del 19-05-2026 & Anton R. Rupi & 2026-05-26 \\[4pt]

    \ghissue{172} & - & Diario di bordo 22-05-2026 & Anton R. Rupi & 2026-05-22 \\[4pt]

    \ghissue{173} & VI-14.3 & Studio e realizzazione automazioni per metriche cruscotto mancanti & Francesco Marcon & 2026-05-26 \\[4pt]

    \ghissue{174} & VI-14.3 & Stesura metriche cruscotto già esistenti & Filippo Idiometri & 2026-05-26 \\[4pt]

	\ghissue{175} & VI-14.4 & Sistemare adr spostando UC1 e UC2 e vanno rinumerati tutti gli uc e i relativi riferimenti in tutto il documento & Armando M. Scarati & 2026-05-26 \\[4pt]

	\ghissue{176} & VI-14.5 & Aggiornare norme di progetto in merito ad associazione ruolo-documento & Sofia De Blasi & 2026-05-26 \\[4pt]

	\ghissue{177} & - & Stesura Sprint 6 & Sofia De Blasi & 2026-05-26 \\[4pt]

	\ghissue{178} & - & Aggiornamento sezione Pianificazione & Sofia De Blasi & 2026-05-26 \\[4pt]
	
\end{longtable}

% --- 6. PROSSIMA RIUNIONE ---
\section{Prossima Riunione}
\begin{itemize}
	\item \textbf{Data:} 2026-05-26
	\item \textbf{Ora:} 14:30
	\item \textbf{OdG preliminare\textsubscript{G}:}
	\begin{itemize}
		\item Revisione e verifica dei task del settimo sprint;
		\item Pianificazione del ottavo sprint e assegnazione dei ruoli.
	\end{itemize}
\end{itemize}